% titel: Blatt 8, A2: Wasser höherer Siedepunkt als H2S
% Teil II: der durchdachte Weg.

\subsection*{Gegeben und gesucht}
Warum siedet H$_2$O (100\,°C) wesentlich höher als H$_2$S (ca. -60\,°C)?
Was würde für eine umgekehrte Reihenfolge sprechen?

\subsection*{Begründung}
a) H$_2$O ist stark polar und bildet ausgeprägte \textbf{Wasserstoffbrücken}
zwischen den Molekülen (H an O). Diese sind weit stärker als die schwachen
Dipol-Dipol- und Dispensionskräfte in H$_2$S, das fast keine H-Brücken bildet
(S ist nur schwach elektronegativ). Man braucht deshalb viel mehr Energie,
um die H$_2$O-Moleküle voneinander zu trennen $\Rightarrow$ hoher Siedepunkt.

\subsection*{Was für eine umgekehrte Reihenfolge spräche}
b) Ohne die H-Brücken-Überhöhung bestimmte allein die \textbf{Molmasse} den
Siedepunkt (grössere Teilchen $\Rightarrow$ stärkere Dispersion). Da H$_2$S
(M$\approx$34) schwerer ist als H$_2$O (M$\approx$18), würde man dann den
höheren Siedepunkt natürlich beim H$_2$S erwarten. Genau dieser
Masseneffekt wird durch die H-Brücken des Wassers umgekehrt.

\begin{ergebnis}
a) H$_2$O bildet wesentlich stärkere Wasserstoffbrücken als H$_2$S.
b) Nach der Molmasse (Dispersion) läge H$_2$S oben --- die H-Brücken drehen
die Reihenfolge um.
\end{ergebnis}

\begin{kontrolle}
Zeigt eine Tabelle O- vs. S-Verbindungen (H$_2$O, H$_2$S, H$_2$Se, H$_2$Te),
so bildet nur H$_2$O den steilen Ausreisser nach oben --- konsistent mit
einzigartig starken H-Brücken.
\end{kontrolle}

\begin{merksatz}
Wenn eine Verbindung im Siede- oder Schmelzpunkt stark aus einer regelmässigen
Reihe herausragt, sucht man einen qualitativen Sonderfaktor
(Wasserstoffbrücken, Ionen, besonders polare Bindungen) und nennt ihn.
\end{merksatz}
